\begin{abstract}
Storage systems already expose the ingredients for autonomous maintenance: device-health observations, checksums and redundancy, scrub and repair primitives, data migration, and workload telemetry. Turning these ingredients into a closed loop is nevertheless risky. Device-health signals are noisy, a false alarm can trigger unnecessary data movement, and repair traffic competes with foreground I/O; most importantly, an automated maintenance action can itself reduce the margin for recovery if it is issued in an unsafe state.

We present \autopilot{}, a safety-first maintenance controller implemented in \argosfs{}. The controller treats health prediction as uncertain evidence rather than an oracle. It keeps persistent per-device evidence across cycles, requires confirmation or critical evidence before proactive drain, applies structural safety gates before mutating storage state, bounds scrub and rebalance to resumable work windows, reduces or pauses background work under foreground pressure, and records action outcomes for later decisions. The controller coordinates four actions---observe, scrub, rebalance, and drain---without introducing a new failure predictor.

We evaluate the design against reactive, periodic, fixed-threshold, and oracle policies using deterministic fault traces, replay of public drive-health traces, real FUSE workloads, and device-mapper fault injection. The final empirical summary is \resulttext{insert the final safety, risk-exposure, and foreground-tail-latency findings.} The implementation and experiment drivers are open and retain seeds, controller decisions, and raw measurements for replay.
\end{abstract}
